Bei der Brown'schen Bewegung beobachtet man unter dem Mikroskop kleine Teilchen (z.B. Fetttröpfchen in Milch oder Rauchteilchen in Luft), die von den unsichtbaren Molekülen ihrer Umgebung unregelmässig angestossen werden.

\begin{itemize}
\item Jedes Teilchen zittert ständig hin und her und ändert dauernd Richtung und Betrag seiner Geschwindigkeit. Die Teilchen sind nie in Ruhe, der Betrag ihrer Geschwindigkeit ist im Mittel also grösser als null: \textbf{Die mittlere Schnelligkeit ist nicht null.}
\item Die Bewegung hat keine Vorzugsrichtung: Die Teilchen wandern im Mittel nicht in eine bestimmte Richtung ab, ihr gemeinsamer Schwerpunkt bleibt an Ort (die Milch fliesst nicht weg). Jede Geschwindigkeitskomponente kommt mit positivem und negativem Vorzeichen gleich häufig vor, die Geschwindigkeitsvektoren heben sich im Mittel auf: \textbf{Die mittlere Geschwindigkeit (als Vektor) ist null.}
\end{itemize}
